\section{Hipervisor Tipo 1}

Un hipervisor Tipo 1 se ejecuta directamente sobre el hardware; uno Tipo 2 depende de un sistema operativo anfitrión. Proxmox sustituyó el sistema operativo de la torre dedicada, por lo que no existe una capa de VirtualBox, VMware Workstation u otro hipervisor entre Proxmox y el equipo físico \cite{proxmox-install}.

\begin{table}[H]
  \centering
  \begin{tabularx}{\linewidth}{>{\bfseries}lYY}
    \toprule
    Criterio & Implementación del laboratorio & Alternativa no utilizada \\
    \midrule
    Tipo & Tipo 1, bare-metal & Tipo 2 sobre Windows \\
    Plataforma & Proxmox VE 9.2 & VirtualBox/VMware \\
    Máquinas & KVM/QEMU 101, 102 y 103 & VM anidadas \\
    Condición de puntos extra & Conservada & Se anularía \\
    \bottomrule
  \end{tabularx}
  \caption{Diferencia práctica entre la implementación elegida y virtualización anidada.}
\end{table}

\evidencia{Versión y recursos del host Proxmox}{01-proxmox/E01-proxmox-version.jpeg}{La consola corresponde a \texttt{root@pve}, muestra la rama Proxmox VE 9.2, almacenamiento local activo y ocho banderas de virtualización detectadas en el CPU físico. Esto documenta el host, no una de las VM.}

\evidencia{Tres máquinas virtuales activas en Proxmox}{01-proxmox/E02-proxmox-tres-vm.png}{La interfaz de Proxmox muestra las VM 101 \texttt{lab5-web}, 102 \texttt{lab5-db} y 103 \texttt{lab5-cliente} con indicadores verdes, uso de CPU, memoria y tiempo de actividad simultáneo.}
